% C5: Hyperparameter Sensitivity Analysis
% Shows that main routing performance is robust to hyperparameter choices.
% Cross-references: Figure 4, A0 (pseudocode), D1 (configuration)

\subsection{Hyperparameter Sensitivity}
\label{appendix:hp_sensitivity}

We verify that the Figure~4 results are not an artifact of hyperparameter
tuning by sweeping the two most impactful hyperparameters: the Corralling
meta-learner learning rate~$\eta$ and the UCB exploration coefficient
starting value~$\alpha_{\text{start}}$.  All runs use the same protocol as
Figure~4 ($\lambda{=}0$, 20 seeds, dev-set burn-in, holdout evaluation).

\subsubsection{Corralling Learning Rate ($\eta$)}

We fix $\alpha_{\text{start}}{=}2.0$ and sweep
$\eta \in \{0.1, 0.3, 0.5, 1.0, 2.0, 5.0\}$.

\begin{table}[h]
\centering
\caption{Sensitivity to Corralling learning rate $\eta$ ($\lambda{=}0$,
$\alpha_{\text{start}}{=}2.0$, 20 seeds).}
\label{tab:eta_sensitivity}
\small
\begin{tabular}{@{}ccc@{}}
\toprule
$\eta$ & \textbf{Reward} & \textbf{95\% CI} \\
\midrule
0.1 & 0.9140 & $\pm$0.0030 \\
0.3 & 0.9161 & $\pm$0.0017 \\
0.5 & 0.9138 & $\pm$0.0040 \\
1.0 (default) & 0.9110 & $\pm$0.0057 \\
2.0 & 0.9135 & $\pm$0.0056 \\
5.0 & 0.9137 & $\pm$0.0053 \\
\bottomrule
\end{tabular}
\end{table}

\paragraph{Finding.}
Performance is remarkably stable across a 50$\times$ range of $\eta$ (0.1 to
5.0), with all rewards falling within 0.9110--0.9161 (a 0.5\% spread).  The
default $\eta{=}1.0$ achieves 0.9110, within CI of the best (0.9161 at
$\eta{=}0.3$).  Higher $\eta$ values produce wider confidence intervals,
consistent with more volatile meta-weight switching, but the mean reward
difference is negligible.

\subsubsection{UCB Exploration Coefficient ($\alpha_{\text{start}}$)}

We fix $\eta{=}1.0$ and sweep
$\alpha_{\text{start}} \in \{0.5, 1.0, 2.0, 4.0\}$, with linear decay to
$\alpha_{\text{end}}{=}0.1$.

\begin{table}[h]
\centering
\caption{Sensitivity to UCB exploration coefficient $\alpha_{\text{start}}$
($\lambda{=}0$, $\eta{=}1.0$, 20 seeds).}
\label{tab:alpha_sensitivity}
\small
\begin{tabular}{@{}ccc@{}}
\toprule
$\alpha_{\text{start}}$ & \textbf{Reward} & \textbf{95\% CI} \\
\midrule
0.5  & 0.9129 & $\pm$0.0049 \\
1.0  & 0.9110 & $\pm$0.0074 \\
2.0 (default)  & 0.9121 & $\pm$0.0095 \\
4.0  & 0.9129 & $\pm$0.0028 \\
\bottomrule
\end{tabular}
\end{table}

\paragraph{Finding.}
Similarly robust: all $\alpha_{\text{start}}$ values in $[0.5, 4.0]$ yield
rewards in 0.9110--0.9129 (0.2\% spread).  All values are within each
other's confidence intervals, confirming that the system is insensitive
to the exploration schedule.  The default $\alpha_{\text{start}}{=}2.0$
is within the plateau.

\subsubsection{Implications}

Both $\eta$ and $\alpha_{\text{start}}$ exhibit a broad plateau:
performance degrades by $<0.5\%$ across order-of-magnitude hyperparameter
changes.  This robustness is a practical strength: practitioners can adopt
the default hyperparameters without task-specific tuning.  The stability
arises from two factors.  First, the Corralling safety mechanism
($\gamma$-mixing ensures all experts receive samples regardless of $\eta$).
Second, LinUCB's ridge regression is inherently regularized, dampening
sensitivity to the exploration schedule.  No hyperparameter tuning was
performed on the holdout set.

\subsubsection{Recommended Defaults}

Based on the sensitivity analysis, we recommend the following defaults for
new deployments:

\begin{itemize}[leftmargin=*,noitemsep]
    \item $\eta{=}1.0$ for \textbf{standard routing} (quality/cost optimization).
          Performance is insensitive to $\eta$ in the $[0.1, 5.0]$ range,
          so this value requires no task-specific tuning.
    \item $\eta{=}0.3$ for \textbf{safety-critical deployments} where
          catastrophic failure detection is needed with low false-positive
          rates (${\leq}10\%$).
    \item $\alpha_{\text{start}}{=}2.0$ with linear decay to $0.1$.
          Sufficient exploration for burn-in sets of $500{+}$ prompts.
          For smaller burn-in sets ($< 200$), consider $\alpha_{\text{start}}
          {=}4.0$ to increase early exploration.
    \item All hyperparameters can be used as-is without holdout tuning.
\end{itemize}

\subsubsection{Catastrophic Failure Detection ($\eta$ Ablation)}
\label{appendix:hp_catastrophic_eta}

For the catastrophic failure experiment (Appendix~\ref{appendix:catastrophic_failure}),
where $\eta$ controls the speed of expert decommissioning, the trade-off is
more nuanced.  We sweep $\eta \in \{0.1, 0.3, 0.5, 1.0, 2.0, 5.0\}$ over
20 seeds on the three-phase catastrophic failure environment (500 steps each).

\begin{table}[h]
\centering
\caption{Learning rate sensitivity for catastrophic failure detection (20
seeds, 500 steps).  Detection = steps after failure onset to decommission;
FP = premature decommissioning in healthy Phase~1.}
\label{tab:eta_catastrophic}
\small
\begin{tabular}{@{}ccccc@{}}
\toprule
$\eta$ & \textbf{Detection (steps)} & \textbf{FP Rate} & \textbf{Recovery Rate} & \textbf{Recovery (steps)} \\
\midrule
0.1  & $41.0 \pm 19.3$ & 0\%   & 0\%  & --- \\
0.3  & $14.3 \pm 11.4$ & 10\%  & 0\%  & --- \\
0.5  & $6.9 \pm 7.4$   & 40\%  & 10\% & $27.5 \pm 7.5$ \\
1.0  & $11.9 \pm 18.9$ & 40\%  & 10\% & $12.0 \pm 2.0$ \\
2.0  & $5.8 \pm 9.5$   & 75\%  & 15\% & $53.3 \pm 60.5$ \\
5.0  & $6.1 \pm 11.3$  & 100\% & 10\% & $92.0 \pm 84.0$ \\
\bottomrule
\end{tabular}
\end{table}

\paragraph{Trade-off.}
In the catastrophic failure setting, $\eta$ has a more pronounced effect,
enabling the system to quickly detect and decommission failing models.
$\eta{=}0.3$ (used in the main paper) achieves $3\times$ faster detection
than $\eta{=}0.1$ (14.3 vs.\ 41.0 steps) while keeping the false-positive
rate at only 10\%.  At higher values ($\eta \geq 0.5$), detection accelerates
further but at the cost of $\geq 40\%$ false positives --- prematurely
decommissioning healthy experts.  The user can select $\eta$ based on their
deployment priority: $\eta{=}0.1$--$0.3$ for safety-critical applications
where false alarms are costly, or $\eta{=}0.5$--$1.0$ for latency-critical
applications where rapid failure detection takes priority.

\subsubsection{Reproducibility}

\begin{itemize}[noitemsep]
    \item Routing task: \texttt{experiments/04\_figure/run\_hyperparameter\_sensitivity.py}
    \item Catastrophic failure: \texttt{experiments/appendix/E\_catastrophic\_failure\_experiment/supplementary/ablation\_learning\_rate\_catastrophic.py}
\end{itemize}
